\documentclass{ctexart}
\usepackage{geometry}
\geometry{a4paper, left=3.08cm, right=3.08cm, top=2.54cm, bottom=2.54cm} % word文档默认页边距
\usepackage[dvipsnames]{xcolor} % 颜色宏包
\usepackage{amsfonts} % 数学字体
\usepackage{amsmath} % 数学公式
% 三线表宏包
\usepackage{booktabs}
\usepackage{makecell}
% 图形宏包
\usepackage{graphicx} 
\usepackage{subcaption}	
\usepackage{float}
% 跨多行多列的表格
\usepackage{multirow}
%\setlength{\baselineskip}{20pt}
%\usepackage{titlesec}
%\title{My Custom LaTeX Template}
%\author{Your Name字}
%\pagestyle{plain} % 没有页眉，页脚只包含一个居中的页码
%代码宏包
\usepackage{listings}
\lstset{
	frame=tb,
	aboveskip=3mm,
	belowskip=3mm,
	showstringspaces=false,
	columns=flexible,
	framerule=1pt,
	rulecolor=\color{gray!35},
	backgroundcolor=\color{gray!5},
	basicstyle={\small\ttfamily},
	numbers=none,
	numberstyle=\tiny\color{gray},
	keywordstyle=\color{blue},
	commentstyle=\color{OliveGreen},
	breaklines=true,
	breakatwhitespace=true,
	tabsize=3,
}
\renewcommand {\lstlistingname}{源程序}



\begin{document}
\pagestyle{empty}
\quad	
\vspace{7ex}
	
\begin{figure}[htb]
	\centering
	\includegraphics[width=0.75\linewidth]{./figure/head}
	\label{fig:head}
\end{figure}

\begin{center} 
	\zihao{-0}\bfseries\heiti{课程设计报告书}\vspace{2ex}\\
	\zihao{2}\bfseries\heiti{题目：}\vspace{3ex}\\
	\zihao{-3}\songti{学院：数学学院}\vspace{1ex}\\
	\zihao{-3}\songti{专业：数学与应用数学}\vspace{1ex}\\
	\zihao{-3}\songti{学生姓名：BuggingCat}\vspace{1ex}\\
	\zihao{-3}\songti{学生学号：2021}\vspace{1ex}\\
	\zihao{-3}\songti{指导老师：}\vspace{1ex}\\
	\zihao{-3}\songti{课程编号：}\vspace{1ex}\\
	\zihao{-3}\songti{课程学分：2}\vspace{1ex}\\
	\zihao{-3}\songti{起始日期：2024.01}
\end{center}

\clearpage

\begin{table}[]
	\centering
	\begin{tabular}{|l|l|}
		\hline
		\begin{tabular}[c]{@{}l@{}}\\ \\ \\ \\ \\ \\ 教\\ 师\\ 评\\ 语\\ \\ \\ \\ \\ \\ \\ \end{tabular} & \begin{tabular}[c]{@{}l@{}}\\ \\ \\ \\ \\ \\ \\ \\ \\ \\ \\ \\ \\ \\ \\ \\ \\ \\ \\ \\ \qquad \qquad \qquad \qquad \qquad \qquad \qquad \qquad \qquad \qquad \qquad \qquad \qquad \qquad 教师签名：\qquad \qquad \qquad\\ \qquad \qquad \qquad \qquad \qquad \qquad \qquad \qquad \qquad \qquad \qquad \qquad \qquad \qquad 日期：\end{tabular} \\ \hline
		\begin{tabular}[c]{@{}l@{}}\\ \\ 成\\ 绩\\ 评\\ 定\\ \\ \\ \end{tabular} &                                                                                                                 \\ \hline
		\begin{tabular}[c]{@{}l@{}}\\ \\ 备 \\ 注\\ \\ \\ \end{tabular}                                                    &                                                                                                                 \\ \hline
	\end{tabular}
\end{table}

\clearpage

\tableofcontents % 目录

\clearpage

\pagestyle{plain}
\setcounter{page}{1} % 页码从中间开始编号

\begin{center} 
% 标题
% 三号宋体加粗，上下空一行
\zihao{-2}\bfseries\heiti{论文标题}
%\vspace{2ex}
\end{center}

% 正文内容
\section{选题背景}

说明本课题应解决的主要问题及应达到的技术要求，简述本设计的指导思想。

引用参考文献\cite{ref1}，引用图片\ref{fig:first}，引用公式\ref{eq1}，引用表格\ref{tab:9}。

\begin{table}[htbp]
	\centering
	\caption{数据集的LPIPS值.}
	\label{tab:9}
	\begin{tabular}{l|cc|cc}
		%\toprule
		& \textbf{a} & \textbf{b} & \textbf{c} & \textbf{d} \\
		\hline
		A   & 0.2324  & 0.5678  & 0.5656    & 0.1256    \\
		B & 0.2123  & 0.3478  & 0.3812    & 0.2421    \\
		C & 0.3456  & 0.3896  & 0.3521    & 0.4231    \\
		D & 0.1234  & 0.2345  & \textbf{0.2372}    & 0.2520    \\
		
		%\bottomrule
	\end{tabular}
\end{table}


\section{方案论证（设计理念）}

说明设计原理（理念）并进行方案选择，阐明为什么要选择这个设计方案以及所采用方案的特点。

包括：重点说明要实现的功能及其要求、系统的安全性、数据的完整性、应用的运行环境及其性能等要求。

\section{过程论述}

重点说明设计是如何实现的，包括：对设计工作的详细表述。要求层次分明、表达确切。

要求：每个图都必须有文字说明，图前说明为什么使用该图、图的主要作用；图后说明图中各成分的作用，和成分之间的交互或图所表达的流程。

\subsection{什么}

字字字字字，吱吱字，字如下：
\begin{description}
	\item[第1步:] 数据预处理/参数设定
	\item[第2步:] 应用XXX离散格式/函数来计算XXX
	\item[第3步:] 利用XXX方法可得XXX
\end{description}

\subsection{什么什么}

吱吱吱吱，字，吱吱字，字字。

\begin{equation}\label{eq1}
	i\hbar\frac{\partial}{\partial t} \Psi(r,t)  = \left[
	-\frac{\hbar^2}{2m}\nabla^2+V(r,t) \right]\Psi(r,t) 
\end{equation}

\begin{itemize}
	\item[1.]不需要把所有变量都放到这个表里面，模型中用到的临时变量可以不放。
	\item[2.]下文中首次出现这些变量时也要进行解释，不然会降低文章的可读性。
\end{itemize}

\begin{table}[h]
	\centering
	表1\quad 什么什么的表\vspace{1ex}
	\begin{tabular}{p{2cm}<{\centering}p{9cm}<{\centering}p{2cm}<{\centering}}
		\toprule
		符号 & 说明 & 单位 \\
		\midrule
		$\alpha$ & 建模花费时间  & 天   \\
		$\mathbb{E}[X]$ & 股票收盘价$X$的数学期望 & 元 \\
		$W_{\text{max}}$ & 某人体重$W$的最大值 & 千克  \\
		$x_{ij}$ & 第$i$辆货车运输第$j$种货物的单位费用&  元/吨 \\
		\bottomrule
	\end{tabular}
\end{table}

吱吱吱吱，字，吱吱字，字字。什么什么的什么是什么什么，什么的什么是什么什么的什么，字字字字字字字，吱吱字字字。

\begin{figure}[htb]
	\centering
	\includegraphics[width=1\linewidth]{./figure/logo}
	\caption{华南理工大学校徽、校名与数学数学院徽}
	\label{fig:logo}
\end{figure}

\section{结果分析}

对研究过程中所获得的主要的数据、现象进行定性或定量分析，得出结论和推论。

引用图片\ref{fig:first}。


\begin{figure}[H]
	\centering
	\begin{subfigure}{0.4\textwidth}
		\includegraphics[width=\textwidth,angle=90]{./figure/Cat}
		\caption{第一幅图，可以设置旋转角度}
		\label{fig:first}
	\end{subfigure}
	\hfill
	\begin{subfigure}{0.4\textwidth}
		\includegraphics[width=1.2\textwidth]{./figure/watermarked}
		\caption{第二幅图}
		\label{fig:second}
	\end{subfigure}
	\hfill
	\begin{subfigure}{0.85\textwidth}
		\includegraphics[width=\textwidth]{./figure/flag}
		\caption{第三幅图}
		\label{fig:third}
	\end{subfigure}
	
	\caption{可以添加多幅图}
	\label{fig:figures}
	
\end{figure}


\section{课程设计总结}

总结可以包括:课程设计过程的收获、遇到的问题，遇到问题解决问题过程的思考、程序调试能力的思考，课程设计实现过程中的收获和体会等。

\begin{thebibliography}{20}
	\bibitem{ref1}郭莉莉,白国君,尹泽成,魏惠芳. “互联网+”背景下沈阳智慧交通系统发展对策建议[A]. 中共沈阳市委、沈阳市人民政府.第十七届沈阳科学学术年会论文集[C].中共沈阳市委、沈阳市人民政府:沈阳市科学技术协会,2020:4.
	\bibitem{ref2}陈香敏,魏伟,吴莹. “文化+人工智能”视阈下文化创意产业融合发展实践及路径研究[A]. 中共沈阳市委、沈阳市人民政府.第十七届沈阳科学学术年会论文集[C].中共沈阳市委、沈阳市人民政府:沈阳市科学技术协会,2020:4.
\end{thebibliography}

\section{源程序}

\begin{lstlisting}[language=matlab,caption={求解问题1的MatLab代码}]
	% 代码！也可以写在这里！
	t = [48.7 53.0 61.3 72.0 48.7 52.0 64.0];
	w = [2000 3000 1000 500 4000 2000 1000];
	n = [8 7 9 6 6 4 8];
	clj = 1020;
	cwj = 40000;
	s = 302.7;
	f = -[t t];
	A = [1 0 0 0 0 0 0 1 0 0 0 0 0 0; ...
	0 1 0 0 0 0 0 0 1 0 0 0 0 0; ...
	0 0 1 0 0 0 0 0 0 1 0 0 0 0; ...
	0 0 0 1 0 0 0 0 0 0 1 0 0 0; ...
	0 0 0 0 1 0 0 0 0 0 0 1 0 0; ...
	0 0 0 0 0 1 0 0 0 0 0 0 1 0; ...
	0 0 0 0 0 0 1 0 0 0 0 0 0 1; ...
	t zeros(1,7); zeros(1,7) t; ...
	w zeros(1,7); zeros(1,7) w; ...
	zeros(1,4) t(5:7) zeros(1,4) t(5:7)];
	b = [n clj clj cwj cwj s]';
	intcon = 1:length(f);
	lb = zeros(length(f),1);
	[x,fval,exitflag,output] = intlinprog(f,intcon,A,b,[],[],lb)
\end{lstlisting}

\section{如果你不要目录，也可以把tableofcontents注释掉}

\end{document}